\chapter{控制流、函数和错误路径}

\section{问题入口：命令行参数解析}

控制流不是只会写 \code{if} 和 \code{for}。真正的训练是把正常路径和错误路径写清楚。假设程序支持：

\begin{lstlisting}[language=bash,caption={命令行需求}]
tool --input data.txt --limit 100
\end{lstlisting}

朴素写法会在 \code{main} 里一边循环一边判断：

\begin{lstlisting}[language=C++,caption={混乱的参数解析}]
for (int i = 1; i < argc; ++i) {
    std::string arg = argv[i];
    if (arg == "--input") {
        input = argv[++i];
    } else if (arg == "--limit") {
        limit = std::stoi(argv[++i]);
    }
}
\end{lstlisting}

这段代码漏了很多边界：\code{--input} 后面没有值会越界，\code{--limit abc} 会抛异常，重复参数如何处理也不明确。

\section{先定义结果}

参数解析的结果应该是一个类型：

\begin{lstlisting}[language=C++,caption={参数解析结果}]
#include <optional>
#include <string>

struct Options {
    std::string input_path;
    int limit = 0;
};

struct ParseError {
    std::string message;
};
\end{lstlisting}

标准库直到 C++23 才有 \code{std::expected}。为了保持 C++20，本章用一个简单结构表达“成功或失败”：

\begin{lstlisting}[language=C++,caption={C++20 下的简单结果类型}]
struct ParseResult {
    std::optional<Options> options;
    std::optional<ParseError> error;
};
\end{lstlisting}

真实项目可以使用成熟的 \code{expected} 实现；教材这里先把控制流讲清楚。

\section{让循环只做推进}

解析函数可以按下标推进。每次遇到需要值的选项，都先检查后面是否还有参数。

\begin{lstlisting}[language=C++,caption={明确错误路径的参数解析}]
#include <cstdlib>
#include <string_view>
#include <vector>

ParseResult parse_options(std::span<char* const> args)
{
    Options options;
    bool has_input = false;
    bool has_limit = false;

    for (std::size_t i = 1; i < args.size(); ++i) {
        const std::string_view arg = args[i];
        if (arg == "--input") {
            if (i + 1 >= args.size()) {
                return {.error = ParseError{"--input requires a value"}};
            }
            options.input_path = args[++i];
            has_input = true;
            continue;
        }
        if (arg == "--limit") {
            if (i + 1 >= args.size()) {
                return {.error = ParseError{"--limit requires a value"}};
            }
            const std::string value = args[++i];
            try {
                options.limit = std::stoi(value);
            } catch (const std::exception&) {
                return {.error = ParseError{"--limit must be an integer"}};
            }
            has_limit = true;
            continue;
        }
        return {.error = ParseError{"unknown option"}};
    }

    if (!has_input) {
        return {.error = ParseError{"missing --input"}};
    }
    if (!has_limit) {
        return {.error = ParseError{"missing --limit"}};
    }
    return {.options = options};
}
\end{lstlisting}

这段代码不追求最短。它追求每个错误路径都能解释。工程代码里，错误消息是接口的一部分。

\section{函数大小和责任边界}

当解析规则继续增多，可以把“取下一个值”抽成辅助函数，把“字符串转整数”抽成辅助函数。拆函数不是为了好看，而是为了隔离可测试的决策点。

\begin{lstlisting}[language=C++,caption={把字符串转整数变成独立函数}]
std::optional<int> parse_int(std::string_view text)
{
    try {
        std::size_t consumed = 0;
        const int value = std::stoi(std::string{text}, &consumed);
        if (consumed != text.size()) {
            return std::nullopt;
        }
        return value;
    } catch (const std::exception&) {
        return std::nullopt;
    }
}
\end{lstlisting}

这里有一个小成本：\code{std::stoi} 接受 \code{std::string}，所以我们从 \code{string_view} 构造了字符串。性能敏感场景可以用 \code{std::from_chars}。入门阶段先把正确性和接口讲清楚。

\section{本章边界检查}

控制流代码要重点检查：

\begin{itemize}
  \item 选项缺值。
  \item 未知选项。
  \item 数字格式错误。
  \item 重复选项是否允许。
  \item 默认值是否真实表达业务需求。
\end{itemize}

\begin{keyidea}
控制流的高级能力不是写复杂分支，而是让每条路径都有清晰含义。正常路径短，错误路径明确，状态推进单调，函数边界可测试。
\end{keyidea}

\section{从缺值反例推到下标推进规则}

参数解析最容易错在 \code{++i}。看这个输入：

\begin{lstlisting}[language=bash,caption={缺少参数值的输入}]
tool --input
\end{lstlisting}

如果代码直接写 \code{input = argv[++i]}，\code{i} 会越过最后一个参数。这个反例推出第一条规则：任何消耗下一个参数的选项，都必须先检查 \code{i + 1 < args.size()}。再看另一个输入：

\begin{lstlisting}[language=bash,caption={未知选项不能静默忽略}]
tool --input data.txt --limt 100
\end{lstlisting}

\code{--limt} 是拼写错误。如果解析器静默忽略未知选项，程序可能用默认限制继续跑，结果比直接失败更危险。这个反例推出第二条规则：未知选项应该成为错误，除非需求明确允许透传。

\section{让数字解析没有半截成功}

\code{std::stoi("12abc")} 会解析出 \code{12}，如果不检查消费长度，就会把坏输入当成好输入。前面 \code{parse_int} 里检查 \code{consumed != text.size()}，就是为了排除这种半截成功。

在 \cpp20 中也可以用 \code{std::from_chars}，它不抛异常，也不受本地化影响：

\begin{lstlisting}[language=C++,caption={用字符转换解析整数}]
#include <charconv>

std::optional<int> parse_int_fast(std::string_view text)
{
    int value = 0;
    const char* begin = text.data();
    const char* end = text.data() + text.size();
    const auto [ptr, ec] = std::from_chars(begin, end, value);
    if (ec != std::errc{} || ptr != end) {
        return std::nullopt;
    }
    return value;
}
\end{lstlisting}

这个版本更适合底层解析。它的接口仍然返回 \code{std::optional<int>}，因为调用者关心的是“是否完整解析成功”，而不是底层用了异常还是错误码。

\section{把主流程写成可读句子}

解析函数写好后，\code{main} 不应该继续塞满细节。入口可以保持很薄：

\begin{lstlisting}[language=C++,caption={入口只负责连接步骤}]
int run(std::span<char* const> args)
{
    const ParseResult parsed = parse_options(args);
    if (parsed.error.has_value()) {
        std::cerr << parsed.error->message << '\n';
        return 2;
    }

    return run_with_options(*parsed.options);
}
\end{lstlisting}

这里有一个隐含要求：\code{ParseResult} 不能同时既有 \code{options} 又有 \code{error}。更严谨的项目会用 \code{expected} 或自定义变体类型表达二选一。教材这里先用简单结构，是为了把控制流看清楚；真正工程里可以进一步收紧类型。

\begin{exercisebox}
给参数解析器补上重复选项规则：\code{--limit 10 --limit 20} 是报错，还是后者覆盖前者？先写出你的规则，再修改 \code{has_limit} 分支，让测试覆盖这两种输入。
\end{exercisebox}
